\section{Discussion}\label{sec:discussion}

The main limitation is coverage rather than the number of source rows. Although the joined modelling data contain many
records, the ranking catalogue uses only 1,530 accepted Portuguese \textit{final} rows because the other sources do not
provide the same safe scientific crop--pathogen links. These rows form 361 complete contexts and 595 context--BCA
interactions, while the ranking evaluation uses only the 109 contexts with at least two observed BCAs. Repeated
authorisation rows should not be read as independent biological observations. The 98.2\% candidate coverage therefore
means that the held-out BCA remained in the catalogue for 107 of these 109 contexts; it does not indicate broad coverage
of new crop--pathogen pairs.

A pathogen support signal accounts for much of the ranking behaviour in the current data. The sources repeat
authorisation records around target organisms, while a missing record does not provide reliable evidence of absence.
The hybrid ranker slightly improves on direct support alone in HitRate@5, MRR, and NDCG@10. Collaborative support
alone performs worse. The improvement over direct support is small: $+0.9$ percentage points in Hit@5, $+0.006$
MRR, and $+0.005$ NDCG@10. No confidence intervals or significance test are available, so these differences should be treated as
descriptive evidence for a possible small benefit from collaboration rather than as a firm improvement.

The screening-model results also need a narrow interpretation. The grouped held-out set keeps related source, microbe, pathogen,
and product rows on the same side of the split, and rule-derived match and review fields are excluded. However, the
screening outcomes are generated from the same authorisation records and related fields that supply the model features. The
perfect scores therefore show separability or reproduction of pipeline outcomes, not generalisation to independently
judged BCA status or biological outcomes. The discarded state means that no BCA signal was detected; it does not
mean biological ineffectiveness.

The ranking evaluation uses contexts with several observed BCAs. It does not measure performance for a new crop-pathogen
pair, a BCA absent from the candidate catalogue, or a future regulatory snapshot. Scarce verified use-level data also
limit comparisons across countries and make the current evaluation dependent on the Portuguese catalogue. Expanding the
catalogue with inferred links could increase apparent coverage, but would replace a missing observation with an assumption
and would weaken the meaning of the evaluation.
The screening outcomes inherit the coverage and regulatory bias of the authorisation sources. Because the screening weight is zero,
the screening model is not part of this ranking evaluation. The experiment uses the authorisation-derived final catalogue
and does not test end-to-end generalisation.
